\chapter{Introduction}

\label{sec:intro}

\comLM{Second-to-last thing to be written. Don't sweat it. Leave it to \textit{Papa}
\begin{itemize}
    \item Describe NPP and notation; motivation with historical data (MANY CITATIONS).
    \item The NPP as a doubly-intractable problem (the prior is intractable)
    \item Mention that we start with approximate methods and then move on to exact stuff
\end{itemize}
}

% Here we talk about informative priors
In statistical modelling with historical data, a central challenge is effectively integrating information from past studies with current data.
This is particularly relevant in clinical trials and medical research, where leveraging historical data can improve the efficiency and robustness of statistical inferences.
A natural approach in this setting is the use of \textit{informative priors}, which incorporate prior knowledge into the analysis, thereby enhancing parameter estimation and decision-making \cite{spiegelhalter2004, gelman2013, ibrahim2015}.

% Here we talk about power priors
Among the various methods for constructing informative priors, \textit{power priors} \cite{ibrahim2000, ibrahim2015} have gained significant attention.
They allow for controlled borrowing of information from historical data by adjusting a power parameter, $\eta$, which typically ranges between 0 and 1.
When $\eta = 0$, no historical information is incorporated, whereas $\eta = 1$ fully integrates the historical data into the current analysis.
This flexibility makes power priors particularly appealing in adaptive clinical trial design and other settings where historical evidence is available, but its relevance is uncertain.

However, when $\eta$ is itself treated as unknown and assigned a prior distribution, the resulting posterior is \textit{doubly intractable}: the normalised power prior contains a normalising constant $c_0(\eta)$ that depends on $\eta$ and cannot be evaluated in closed form for general likelihoods, preventing direct use of Metropolis--Hastings.
Current approaches rely on \textit{approximate methods} \cite{carvalho2021} that lack theoretical guarantees and explicit convergence bounds.
This work addresses this gap by developing exact Markov chain Monte Carlo algorithms \cite{murray2012, goncalves2017, vats2021} that avoid evaluating the intractable normalising constant ratio altogether, providing a rigorous computational baseline for inference under the normalised power prior.

